
\documentclass[tikz,border=2pt]{standalone}
\usepackage[utf8]{inputenc}
\usepackage[T1]{fontenc}
\usepackage[french]{babel}
\usepackage{amsmath,amssymb}
\usepackage[table]{xcolor}
\usepackage{siunitx}
\usepackage[version=4]{mhchem}
\usepackage{tabularx,booktabs,array}
\usepackage{tikz}
\usetikzlibrary{arrows.meta,positioning,calc,fit,decorations.pathmorphing,patterns,shapes.geometric}
\usepackage{pgfplots}
\pgfplotsset{compat=1.18}
\definecolor{primary}{RGB}{14,116,144}
\definecolor{accent}{RGB}{16,185,129}
\definecolor{dark}{RGB}{15,23,42}
\definecolor{bglight}{RGB}{241,245,249}
\definecolor{borderlight}{RGB}{203,213,225}
\definecolor{examplepink}{RGB}{236,72,153}
\definecolor{remarkblue}{RGB}{14,165,233}
\definecolor{notepurple}{RGB}{99,102,241}
\definecolor{synthgreen}{RGB}{5,150,105}
\definecolor{synthorange}{RGB}{249,115,22}
\definecolor{synthviolet}{RGB}{79,70,229}
\definecolor{primary}{RGB}{14,116,144}
\definecolor{accent}{RGB}{16,185,129}
\definecolor{dark}{RGB}{15,23,42}
\definecolor{bglight}{RGB}{241,245,249}
\definecolor{examplepink}{RGB}{236,72,153}
\definecolor{remarkblue}{RGB}{14,165,233}
\definecolor{notepurple}{RGB}{99,102,241}
\definecolor{synthgreen}{RGB}{5,150,105}
\definecolor{synthorange}{RGB}{249,115,22}
\definecolor{synthviolet}{RGB}{79,70,229}
\begin{document}
\begin{tikzpicture}[
 scale=0.85,
 transform shape,
 >=Stealth,
 thick,
 lens/.style={line width=1.2pt, primary},
 ray/.style={line width=0.7pt, red!70!black, ->},
 ray2/.style={line width=0.7pt, blue!70!black, ->},
 label/.style={font=\small}
]

% Axe optique
\draw[->] (-1,0) -- (14.5,0) node[right]{\small axe optique};

% Objectif L1
\draw[lens] (2,-2.5) -- (2,2.5);
\draw[lens, fill=primary!10] (1.85,-2.5) -- (2.15,-2.5) -- (2.15,2.5) -- (1.85,2.5) -- cycle;
\draw[primary, ->, line width=0.8pt] (2,2.2) -- (2.25,2.5);
\draw[primary, ->, line width=0.8pt] (2,-2.2) -- (1.75,-2.5);
\node[above, primary, font=\bfseries] at (2,2.7) {$L_1$ (Objectif)};
\node[below, label] at (2,-2.8) {$O_1$};
\fill (2,0) circle (2pt);

% Foyers de L1
\fill[red!70!black] (8,0) circle (2.5pt);
\node[below, label, red!70!black] at (8,-0.3) {$F'_1 = F_2$};

% Oculaire L2
\draw[lens] (10,-1.8) -- (10,1.8);
\draw[lens, fill=primary!10] (9.85,-1.8) -- (10.15,-1.8) -- (10.15,1.8) -- (9.85,1.8) -- cycle;
\draw[primary, ->, line width=0.8pt] (10,1.5) -- (10.25,1.8);
\draw[primary, ->, line width=0.8pt] (10,-1.5) -- (9.75,-1.8);
\node[above, primary, font=\bfseries] at (10,2.0) {$L_2$ (Oculaire)};
\node[below, label] at (10,-2.1) {$O_2$};
\fill (10,0) circle (2pt);

% Distances focales
\draw[<->, synthorange, line width=0.8pt] (2,-3.4) -- (8,-3.4);
\node[below, synthorange, font=\small\bfseries] at (5,-3.4) {$f'_1$};

\draw[<->, synthviolet, line width=0.8pt] (8,-3.4) -- (10,-3.4);
\node[below, synthviolet, font=\small\bfseries] at (9,-3.4) {$f'_2$};

\draw[<->, dark, line width=0.9pt] (2,-4.3) -- (10,-4.3);
\node[below, dark, font=\small\bfseries] at (6,-4.3) {$L = f'_1 + f'_2$};

% Rayons venant de l'infini (parallèles inclinés)
\draw[ray] (-0.5,1.8) -- (2,1.8);
\draw[ray] (-0.5,0.6) -- (2,0.6);

% Rayons après L1 convergent vers F'1 avec décalage B1
\draw[ray] (2,1.8) -- (8,-0.9);
\draw[ray] (2,0.6) -- (8,-0.9);

% Image intermédiaire B1
\fill[examplepink] (8,-0.9) circle (2.5pt);
\node[right, examplepink, font=\small\bfseries] at (8.15,-1.2) {$B_1$};

% Rayons de B1 vers L2
\draw[ray2] (8,-0.9) -- (10,-0.9);
\draw[ray2] (8,-0.9) -- (10,0.5);

% Rayons parallèles après L2 (vers l'oeil)
\draw[ray2] (10,-0.9) -- (13,-1.95);
\draw[ray2] (10,0.5) -- (13,-0.55);
\draw[ray2, dashed] (10,0) -- (13,-1.05);
\node[blue!70!black, font=\scriptsize, align=center] at (12.15,0.75)
{faisceau\\\'{e}mergent parall\`{e}le};

% Angle theta
\draw[dark, thick] (0.3,0) arc[start angle=0, end angle=12, radius=1.8];
\node[right, dark, font=\small] at (1.0,0.4) {$\theta$};

% Angle theta'
\draw[dark, thick] (11.1,0) arc[start angle=0, end angle=-19.3, radius=1.35];
\node[right, dark, font=\small] at (11.7,-0.35) {$\theta'$};

% Oeil
\node[font=\Large] at (13.5,-1.2) {$\odot$};
\node[below, font=\small] at (13.5,-1.6) {\oe il};

\end{tikzpicture}
\end{document}
